\section{Replicacion y tolerancia a fallos}

El banco replica su estado desde un nodo \emph{lider} hacia uno o varios nodos
\emph{replica} mediante un log de transacciones ordenado por numero de
secuencia y un transporte TCP directo lider a replica (no se usa Pub/Sub). El
lider acepta conexiones de replica, las pone al dia (\emph{catch-up}) desde la
secuencia que cada una reporta y despues les empuja en vivo cada transaccion
nueva. Si una replica se cae y vuelve, reanuda exactamente donde se quedo: si
proceso hasta la secuencia 4500, al reconectar pide y recibe la 4501 en
adelante. Asi, apagar y encender replicas no pierde ni duplica transacciones.

\subsection{Log por numero de secuencia}

La pieza central es \codigo{com/escom/banco/data/RegistroTransacciones.java}, un
log de transacciones indexado por su numero de secuencia. Cada
\codigo{com/escom/banco/model/Transaccion.java} es un \codigo{record} con
\codigo{seq}, \codigo{origenId}, \codigo{destinoId} y \codigo{montoCentavos}; el
campo \codigo{seq} fija un orden total. El log se respalda en un
\codigo{ConcurrentSkipListMap}, que mantiene las entradas ordenadas por clave
aunque varios hilos agreguen transacciones de forma concurrente.

\begin{lstlisting}[language=Java,caption={RegistroTransacciones: append y lectura ordenada}]
private final ConcurrentSkipListMap<Long, Transaccion> log = new ConcurrentSkipListMap<>();

public void agregar(Transaccion t) {
    log.put(t.seq(), t);
}

/** Transacciones con seq estrictamente mayor a la dada, en orden de seq. */
public List<Transaccion> desde(long seq) {
    return new ArrayList<>(log.tailMap(seq, false).values());
}
\end{lstlisting}

El metodo \codigo{desde(N)} devuelve todas las transacciones con secuencia
estrictamente mayor a \codigo{N}, ya ordenadas. Esta operacion es la base tanto
del catch-up de una replica que se reincorpora como del streaming en vivo: en
ambos casos el orden lo garantiza el log, no el instante en que se hizo el
\codigo{agregar}.

\subsection{El lider: catch-up y streaming TCP}

El lado lider es \codigo{com/escom/banco/replicacion/ServidorReplicacion.java}.
Abre un \codigo{ServerSocket} en un hilo \emph{daemon} de aceptacion y, por cada
replica que se conecta, lanza un hilo que la atiende. El protocolo de saludo es
una sola linea \codigo{RESUME N}: la replica declara la ultima secuencia que ya
tiene aplicada, y a partir de ahi el lider le envia todo lo que falta y luego lo
nuevo, sirviendose siempre de \codigo{repo.registro().desde(enviado)}.

\begin{lstlisting}[language=Java,caption={ServidorReplicacion: catch-up desde N+1 y streaming en vivo}]
long enviado = leerResume(in.readLine());   // "RESUME N"
while (activo && !sock.isClosed()) {
    List<Transaccion> nuevas = repo.registro().desde(enviado);
    if (nuevas.isEmpty()) {
        Thread.sleep(1);
        continue;
    }
    for (Transaccion tx : nuevas) {
        out.write(CodecTransaccion.aLinea(tx));
        out.write("\n");
        enviado = tx.seq();
    }
    out.flush();
}
\end{lstlisting}

Tras el catch-up (todo lo posterior a \codigo{N}), el mismo bucle convierte el
envio en streaming continuo: cuando no hay nada nuevo duerme un milisegundo y
reintenta; cuando aparecen transacciones recien agregadas al log, las empuja en
orden y avanza el cursor \codigo{enviado} a la ultima \codigo{seq} servida. El
parseo del saludo lo hace \codigo{leerResume}, que ante una linea ausente o mal
formada degrada a \codigo{0} (la replica recibe el log completo).

\begin{nota}
El lider no recuerda por replica cuanto le envio: el cursor \codigo{enviado}
vive en el hilo de esa conexion y se reconstruye en cada reconexion a partir del
\codigo{RESUME N} que manda la replica. El estado de avance lo posee la replica.
\end{nota}

\subsection{La replica: aplicar y reconectar}

El lado replica es \codigo{com/escom/banco/replicacion/ClienteReplicacion.java}.
Corre en su propio hilo daemon: se conecta al lider, envia
\codigo{RESUME} con \codigo{repo.secuenciaActual()} y luego lee linea por linea,
aplicando cada transaccion con \codigo{repo.aplicarReplica(...)}. Si la conexion
se cae, el bucle externo reintenta tras una breve espera y reanuda desde la
secuencia actual de la replica.

\begin{lstlisting}[language=Java,caption={ClienteReplicacion: RESUME, aplicar en vivo y reconexion}]
try (Socket s = new Socket(hostLider, puerto);
     /* ...streams UTF-8... */ ) {

    out.write("RESUME " + repo.secuenciaActual() + "\n");
    out.flush();

    String linea;
    while (activo && (linea = in.readLine()) != null) {
        repo.aplicarReplica(CodecTransaccion.deLinea(linea));
    }
} catch (IOException e) {
    // lider no disponible: se reintenta abajo
}
if (!activo) return;
try { Thread.sleep(500); } catch (InterruptedException e) { return; }
\end{lstlisting}

La serializacion en el cable es responsabilidad de
\codigo{com/escom/banco/replicacion/CodecTransaccion.java}, que codifica cada
transaccion como una linea ASCII de cuatro campos separados por comas
(\codigo{seq,origenId,destinoId,montoCentavos}) con \codigo{aLinea} y la
reconstruye con \codigo{deLinea}. Al aplicar, \codigo{aplicarReplica} mueve el
monto entre cuentas, agrega la transaccion al log local y avanza la secuencia a
\codigo{t.seq()}, de modo que el siguiente \codigo{RESUME} refleje el avance.

\subsection{Rol por entorno y resincronizacion ante fallos}

El rol de cada nodo se decide por variables de entorno en
\codigo{com/escom/banco/server/MainServer.java}: sin \codigo{BANCO\_LIDER\_HOST}
el nodo es LIDER y levanta el \codigo{ServidorReplicacion}; con
\codigo{BANCO\_LIDER\_HOST} es REPLICA y arranca el \codigo{ClienteReplicacion}
apuntando a ese host. El puerto de replicacion sale de
\codigo{BANCO\_REPLICACION\_PUERTO} (por omision \codigo{9090}).

\begin{lstlisting}[language=Java,caption={MainServer: seleccion de rol por entorno}]
if (esLider()) {
    new ServidorReplicacion(CuentaRepository.get(), puertoRepl).iniciar();
    System.out.println("Rol: LIDER - replicacion TCP en puerto " + puertoRepl);
} else {
    String liderHost = System.getenv("BANCO_LIDER_HOST");
    new ClienteReplicacion(CuentaRepository.get(), liderHost, puertoRepl).iniciar();
    System.out.println("Rol: REPLICA - siguiendo a " + liderHost + ":" + puertoRepl);
}
\end{lstlisting}

La tolerancia a fallos emerge de combinar las tres piezas anteriores. Si se
apaga una replica, el lider solo ve cerrarse su socket y el hilo que la atendia
termina; el resto del sistema sigue. Cuando esa replica se enciende de nuevo,
reconecta, manda \codigo{RESUME 4500} y el lider responde con la secuencia 4501
en adelante (catch-up) antes de retomar el streaming en vivo. Como el avance se
deriva siempre del numero de secuencia, no hay huecos ni duplicados.

\begin{advertencia}
La resincronizacion confia en que la secuencia de la replica nunca supere a la
del lider. \codigo{aplicarReplica} es para transacciones ya validadas por el
lider: una replica nunca origina transacciones propias, solo las aplica en el
orden recibido.
\end{advertencia}
